\documentclass[a4paper,11pt,oneside,openright,final,mono]{ufopthesis}

\usepackage[portuguese,brazil]{babel} 
\usepackage[utf8]{inputenc}
\usepackage[T1]{fontenc}
\usepackage{natbib} % [longnamesfirst]
\usepackage{rotating}
\usepackage{a4wide}
\usepackage{graphicx}
\usepackage{amssymb,amscd,eucal,amsmath}
\usepackage{multirow}
\usepackage{subfig}
\usepackage{url}
\usepackage{indentfirst}
\usepackage{amsthm}
\usepackage{amsmath}
\usepackage{amssymb}
%\usepackage[subfig]{tocloft}
%\usepackage[dvips,a4paper,english,bookmarks=true,bookmarksnumbered=true,linktocpage]{hyperref}

%\usepackage[ruled, algochapter, vlined,algo2e]{algorithm2e}
%\renewcommand{\listalgorithmcfname}{Lista de Algoritmos}%
%\renewcommand{\algorithmcfname}{Algoritmo}%


%% vilar 
%\let\insertchapterspaceold=\insertchapterspace
%\renewcommand{\insertchapterspace}{%
%  \insertchapterspaceold
%  \addtocontents{loa}{\protect\addvspace{10pt}}%
%}
%\usepackage{listings} 
%\usepackage{listings}
%\lstset{
%  language=Pascal, % C
%  basicstyle=\fontsize{10}{12}\selectfont,
%  captionpos=b,
%  mathescape=true,
%  escapechar=@,
%  extendedchars=true,
%  showstringspaces=false,
%  columns=fixed,
%  basewidth=0.49em,
%  frame=single,
%  framesep=2mm,
%  xleftmargin=2mm,
%  xrightmargin=2mm,
%  framerule=0.5pt,
%  tabsize=2,
%}
\graphicspath{{figure/}}

% List of files to be processed
\includeonly{
  introduction,
}
\theoremstyle{definition}
\newtheorem{rf}{Requisito Funcional}
\newtheorem{rnf}{Requisito Não Funcional}
\newtheorem{rn}{Regra de Negócio}

\newenvironment{mydescription}%
{ \begin{description}%
%	{$\bullet$}%
	{\setlength{\labelwidth}{30pt}%
	 \setlength{\leftmargin}{35pt}%
	 \setlength{\itemsep}{\parsep}}}%
{ \end{description} }

\sloppy


\begin{document}



% An example of the Portuguese structure follows:
\ufopthesis{
    author={Nome do Aluno},
    university={Universidade Federal de Ouro Preto},
    course={Curso},
    portuguesetitle={Titulo em Português},
    portuguesebrokentitle={Titulo em Português},
    portuguesecourse={Curso em português},
    portugueseshortcourse={Curso encurtado},    
    address={Endereço},
    date={2012-03},
    logo={brasao},
    advisor={Nome Orientador na folha de aprovação}{Dr.}{Universidade Federal de Ouro Preto},
    advisorT={Nome Orientador na Capa},
    coadvisor={Coorientador na folha de aprovação}{Msc.}{Universidade Federal de Ouro Preto},
    coadvisorT={Coorientador na capa},
    member={Nome Membro de Banca}{Dr.}{Universidade Federal de Ouro Preto},
    member={Nome Mebro da Banca}{Dr.}{Universidade Federal de Ouro Preto},
    portugueseabstract={resumo.tex},
    abstract={abstract.tex},
    dedication={dedicatoria.tex},
    ack={agradecimentos.tex},
    }
\input{introduction}

\input{related}

\input{development}

\input{experiments}

\input{conclusion}

\input{appendix}

%\cite{androidDevelopers}

\ufopbibliography{mono}

\end{document}
